% !TeX root = main.tex
\documentclass[a4paper, 12pt]{extreport}

\usepackage[T1,T2A]{fontenc}
\usepackage[utf8]{inputenc}
\usepackage[english,russian]{babel}
\usepackage{amssymb,amsfonts,amsmath,mathtext,cite,enumerate,float}
\usepackage{pgfplots}
\usepackage{graphicx}
\usepackage{tocloft}
\usepackage{listings}
\usepackage{caption}
\usepackage{tempora}
\usepackage{titlesec}
\usepackage{setspace}
\usepackage{geometry}
\usepackage{indentfirst}
\usepackage{pdfpages}
\usepackage{enumerate,letltxmacro}
\usepackage{threeparttable}
\usepackage{hyperref}
\usepackage{flafter}
\usepackage{enumitem}
\usepackage{multirow}
\usepackage{longtable}
\usepackage{array}
\usepackage{ragged2e}

\usepackage[figure,table]{totalcount}
\usepackage{lastpage}

\usepackage{tabularx}
\newcolumntype{Y}{>{\centering\arraybackslash}X}

\usepackage{chngcntr}
\counterwithout{table}{chapter}

\setlist{nosep}

\newcommand{\ssr}[1]{\begin{center}
		\LARGE\bfseries{#1}
	\end{center} \addcontentsline{toc}{chapter}{#1}}

\makeatletter
\renewcommand\LARGE{\@setfontsize\LARGE{22pt}{20}}
\renewcommand\Large{\@setfontsize\Large{20pt}{20}}
\renewcommand\large{\@setfontsize\large{16pt}{20}}
\makeatother
\RequirePackage{titlesec}
\titleformat{\chapter}[block]{\hspace{\parindent}\large\bfseries}{\thechapter}{0.5em}{\large\bfseries\raggedright}
\titleformat{name=\chapter,numberless}[block]{\hspace{\parindent}}{}{0pt}{\large\bfseries\centering}
\titleformat{\section}[block]{\hspace{\parindent}\large\bfseries}{\thesection}{0.5em}{\large\bfseries\raggedright}
\titleformat{\subsection}[block]{\hspace{\parindent}\large\bfseries}{\thesubsection}{0.5em}{\large\bfseries\raggedright}
\titleformat{\subsubsection}[block]{\hspace{\parindent}\large\bfseries}{\thesubsection}{0.5em}{\large\bfseries\raggedright}
\titlespacing{\chapter}{12.5mm}{-22pt}{10pt}
\titlespacing{\section}{12.5mm}{10pt}{10pt}
\titlespacing{\subsection}{12.5mm}{10pt}{10pt}
\titlespacing{\subsubsection}{12.5mm}{10pt}{10pt}

\makeatletter
\renewcommand{\@biblabel}[1]{#1.}
\makeatother

\geometry{left=30mm}
\geometry{right=10mm}
\geometry{top=20mm}
\geometry{bottom=30mm}

\renewcommand{\theenumi}{\arabic{enumi}}
\renewcommand{\labelenumi}{\arabic{enumi}\text{)}}
\renewcommand{\theenumii}{.\arabic{enumii}}
\renewcommand{\labelenumii}{\asbuk{enumii}\text{)}}
\renewcommand{\theenumiii}{.\arabic{enumiii}}
\renewcommand{\labelenumiii}{\arabic{enumi}.\arabic{enumii}.\arabic{enumiii}.}

\renewcommand{\cftchapleader}{\cftdotfill{\cftdotsep}}

\addto\captionsrussian{\renewcommand{\figurename}{Рисунок}}
\DeclareCaptionLabelSeparator{dash}{~---~}
\captionsetup{labelsep=dash}

\captionsetup[figure]{justification=centering,labelsep=dash}
\captionsetup[table]{labelsep=dash,justification=raggedright,singlelinecheck=off}

\graphicspath{{images/}}%путь к рисункам

\newcommand{\floor}[1]{\lfloor #1 \rfloor}

\lstset{
	language=Python,
	basicstyle=\scriptsize\ttfamily,
	commentstyle=\color{green},
	keywordstyle=\color{blue},
	stringstyle=\color{red},
	numbers=left,
	numberstyle=\tiny,
	stepnumber=1,
	numbersep=5pt,
	showstringspaces=false,
	breaklines=true,
	frame=lines,
	captionpos=b,
	tabsize=4,
	extendedchars=true,
	inputencoding=utf8,
	keepspaces=true
}

\lstdefinestyle{cpp_listing}{
	language=C++,
	basicstyle=\scriptsize\ttfamily,
	numbers=left,
	numberstyle=\tiny,
	breaklines=true,
	frame=lines
}

\setlist{nosep}

\pgfplotsset{width=0.85\linewidth, height=0.5\columnwidth}

\linespread{1.3}

\parindent=1.25cm

\def\labelitemi{---}
\setlist[itemize]{leftmargin=1.25cm, itemindent=0.65cm}
\setlist[enumerate]{leftmargin=1.25cm, itemindent=0.55cm}

\newcommand{\specialcell}[2][c]{%
	\begin{tabular}[#1]{@{}c@{}}#2\end{tabular}}

\frenchspacing
\usepackage{silence}
\WarningFilter{latexfont}{Font shape}
\WarningFilter{latexfont}{Some font shapes}
\WarningFilter{caption}{Unused \captionsetup}
\pgfplotsset{compat=1.18}

\begin{document}

\chapter{Аналитическая часть}

\section{Постановка задачи}

Задача поиска элемента в массиве сводится к поиску индекса 
заданного значения $x$ среди $N\ge1$ элементов массива. Если 
искомый элемент отсутствует в массиве, алгоритм должен 
вернуть -1.

В данной работе необходимо реализовать четыре алгоритма, 
решающих поставленную задачу:
\begin{enumerate}
    \item линейный поиск;
    \item итеративный бинарный поиск;
    \item бинарный поиск со степенным сужением интервала;
    \item рекурсивный бинарный поиск.
\end{enumerate}

\section{Линейный поиск}

Линейный поиск работает с исходным массивом и не требует 
его предварительного упорядочивания: необходимо последовательно 
перебрать элементы массива и сравнить каждый с искомым значением $x$. 
Перебор происходит тех пор, пока не будет найдено совпадение. 
Если в результате просмотра всех элементов массива совпадение не 
было найдено, алгоритм возвращает -1.

В лучшем случае, если искомый элемент располагается в самом 
начале массива, выполняется одно сравнение. В худшем случае, 
если элемент расположен в его конце или отсутствует, перебор 
выполняется до границы массива и требует до  шагов. Временная 
сложность алгоритма равна $O(N)$. Так как дополнительная память 
не расходуется, емкостная сложность составляет $0(1)$.

\section{Итеративный бинарный поиск}

Итеративный бинарный поиск применяется К отсортированному массиву. 
Алгоритм работает на замкнутом отрезке, обе границы которого 
включаются в рассматриваемый интервал. На каждом шаге вычисляется 
индекс среднего элемента и выполняется его сравнение с искомым 
элементом $x$ до тех пор, пока рассматриваемый отрезок между левой 
и правой границей не вырожден и содержит хотя бы один элемент. 
Если значение среднего элемента совпадает с искомым, поиск 
завершается возвратом найденного индекса. Если элемент меньше 
искомого, левая граница сдвигается вправо, а если больше, то правая 
граница смещается влево. Если текущий отрезок пуст, но искомый 
элемент не был найден, алгоритм возвращает -1.

В лучшем случае искомый элемент оказывается в середине массива, 
то есть его поиск требует всего одно сравнение. В худшем случае 
на каждой итерации выполняется два сравнения с элементами массива, 
поэтому общее число сравнений оценивается как $2~log_{2}(N+1)$. 
Временная сложность алгоритма равна $O(log~N)$. Так как 
дополнительная память не расходуется, емкостная сложность 
составляет $0(1)$.

\section{Бинарный поиск со степенным сужением интервала}

Бинарный поиск со степенным сужением интервала работает на 
отсортированном массиве и использует побитовые шаги, кратные 
степеням двойки. Сначала определяется наибольшая степень двойки, 
не превышающая длину массива: если длина уже является степенью 
двойки и элемент на этой позиции меньше искомого, область поиска 
перестраивается с учетом оставшейся части массива и ближайшей 
сверху степени двойки. Далее в цикле размер шага уменьшается вдвое, 
пока не дойдет до нуля, а опорная позиция сдвигается вперед 
каждый раз, когда проверяемый элемент меньше искомого. 
После завершения цикла итоговый индекс уточняется по результату 
финального сравнения на текущей позиции.

Временная сложность алгоритма составляет $O(log~N)$. Так как 
дополнительная память под структуры данных не расходуется, 
емкостная сложность составляет $0(1)$.

\section{Рекурсивный бинарный поиск}

Рекурсивный бинарный поиск работает на отсортированном массиве. 
На каждом шаге рассматривается диапазон между левой и правой 
границами. Если в текущем диапазоне не осталось элементов для 
сравнения, поиск завершается возвратом -1, иначе вычисляется 
средний элемент отрезка и сравнивается с искомым значением $x$: при 
совпадении возвращается его индекс, иначе выполняется рекурсивный 
вызов для левой или правой части интервала.

В лучшем случае искомый элемент находится в середине 
массива и находится при первом сравнении. В худшем случае 
число сравнений пропорционально $2~log_{2}(N+1)$. Временная 
сложность алгоритма составляет $O(log~N)$. Так как память 
стека вызовов расходуется под каждый шаг рекурсии, емкостная 
сложность равна $O(log~N)$

\section*{Вывод}

В аналитической части выполнена постановка задачи поиска 
элемента в целочисленном массиве, а также разобраны четыре 
алгоритма, решающих данную задачу: линейный, итеративный 
бинарный поиск, бинарный поиск со степенным сужением 
интервала и рекурсивный бинарный поиск.

\end{document}